% Build from doc/: latexmk -pdf -outdir=build/apresentacao roteiro_apresentacao.tex
\documentclass[11pt,a4paper]{article}
\usepackage[T1]{fontenc}
\usepackage[utf8]{inputenc}
\usepackage[brazilian]{babel}
\usepackage{lmodern}
\usepackage[margin=2.2cm]{geometry}
\usepackage{booktabs}
\usepackage{longtable}
\usepackage{enumitem}
\usepackage{needspace}
\usepackage{hyperref}
\usepackage{xcolor}
\definecolor{accent}{HTML}{007F79}
\hypersetup{colorlinks=true,linkcolor=accent,urlcolor=accent,
  pdftitle={Roteiro de defesa de TCC},
  pdfauthor={Kleber Daniel Mattos Viana}}
\setlength{\parindent}{0pt}
\setlength{\parskip}{6pt}
\setlist{nosep}
\newcommand{\slide}[3]{\Needspace{8\baselineskip}
  \subsection*{Slide #1 --- #2}
  \textbf{Tempo: #3.}\par}
\newcommand{\cue}[1]{\par{\small\color{accent}\textbf{Condução:} #1}\par}

\begin{document}
\begin{center}
  {\LARGE\bfseries Roteiro de defesa de TCC\par}
  \vspace{3mm}
  {\large Kleber Daniel Mattos Viana\par}
  Ciência da Computação --- UNIFESO --- 2026
\end{center}

\textbf{Tema:} Comparação de Métodos de Clusterização na Análise de Máximos
de Tensão em Simulações de Energização de Linhas de Transmissão no ATP.

\textbf{Uso:} fala sugerida para os 20 slides de
\texttt{apresentacao.tex}, voltada a uma banca de Computação sem formação
elétrica. A sequência liga problema, objetivos, dados, implementação, análise probabilística, comparação dos métodos, contribuições e conclusão.
Os tempos incluem explicação das imagens, pausas e transições. São metas
de ensaio; a leitura isolada não garante a duração. Perguntas da banca
ficam fora dos \textbf{20 minutos}.

\section*{Distribuição dos 20 minutos}
\begin{longtable}{r p{9cm} r r}
\toprule
Slide & Assunto & Duração & Acumulado \\
\midrule
1 & Abertura & 0:30 & 0:30 \\
2 & Contextualização: energização de linhas de transmissão & 0:45 & 1:15 \\
3 & Simulação elétrica e problema computacional & 0:50 & 2:05 \\
4 & Objetivos específicos & 0:45 & 2:50 \\
5 & Estudo de caso e pontos de observação & 0:50 & 3:40 \\
6 & Cenários de simulação e composição da base & 0:45 & 4:25 \\
7 & Normalização e interpretação dos valores de tensão & 0:50 & 5:15 \\
8 & Tecnologias e fluxo de processamento & 0:55 & 6:10 \\
9 & Excedência e discrepância: objetivos distintos & 0:50 & 7:00 \\
10 & Distribuição dos máximos e adequação do modelo normal & 1:25 & 8:25 \\
11 & Probabilidade de excedência e incerteza amostral & 1:00 & 9:25 \\
12 & Métodos de clusterização avaliados & 1:35 & 11:00 \\
13 & Agrupamentos: K-Means & 1:10 & 12:10 \\
14 & Agrupamentos: DBSCAN & 1:10 & 13:20 \\
15 & Agrupamentos: Hierárquico & 1:05 & 14:25 \\
16 & Métricas principais: Silhouette e Davies--Bouldin & 1:20 & 15:45 \\
17 & Comparação dos índices em recortes específicos & 1:20 & 17:05 \\
18 & Tamanho amostral e abrangência da comparação & 1:05 & 18:10 \\
19 & Contribuições computacionais e delimitação do escopo & 1:05 & 19:15 \\
20 & Conclusões e continuidade & 0:45 & 20:00 \\
\bottomrule
\end{longtable}

\section*{Fala sugerida}

\slide{1}{Abertura}{0:00--0:30}
Bom dia. Sou Kleber Viana. Este trabalho trata do processamento de resultados
de simulações elétricas para encontrar valores que merecem investigação.
Foi orientado pela professora Andressa Ruback e coorientado pelo professor
Bráulio Chuco. Vou começar pelo contexto elétrico necessário e,
depois, mostrar a solução computacional e seus resultados.
\cue{Olhar para a banca. Não ler o título inteiro.}


\slide{2}{Contextualização: energização de linhas de transmissão}{0:30--1:15}
Uma linha de transmissão transporta energia entre subestações. Energizar significa conectá-la à rede pelo fechamento de uma chave. Logo após essa conexão, a tensão pode oscilar e superar temporariamente sua referência: isso é uma sobretensão. O isolamento precisa suportar esses esforços. Esse é o contexto elétrico; a contribuição deste trabalho está no processamento dos resultados de simulação. Para estudar essas manobras repetidamente, usamos um simulador.
\cue{Explicar o ponto central e fazer a transição ao próximo slide.}


\slide{3}{Simulação elétrica e problema computacional}{1:15--2:05}
Ensaios reais envolvem custos e restrições operacionais. O ATP simula os transitórios; o ATPDraw permite representar o circuito. Variando os instantes de fechamento das chaves, obtemos muitas execuções, registradas em arquivos LIS. O problema computacional é extrair os máximos, conferir sua integridade e comparar registros sem perder a origem. Um valor diferente pode ter várias causas. A rotina seleciona casos para investigação, sem confirmar sozinha erro numérico. Esse problema orienta os objetivos.
\cue{Explicar o ponto central e fazer a transição ao próximo slide.}


\slide{4}{Objetivos específicos}{2:05--2:50}
Os objetivos são extrair e organizar os máximos, verificar a integridade, normalizar e descrever os dados. Depois, estimar excedências com intervalos de confiança e comparar agrupamentos e sinalizações dos métodos. Silhouette e Davies--Bouldin avaliam a geometria das partições. A entrega integra essas etapas em uma rotina reprodutível. Antes de explicar os algoritmos, precisamos definir quais dados entram nessa rotina.
\cue{Explicar o ponto central e fazer a transição ao próximo slide.}


\slide{5}{Estudo de caso e pontos de observação}{2:50--3:40}
O esquema mostra três posições: terminal de manobra, em zero; meio da linha, em 40 quilômetros; e terminal oposto, em 80 quilômetros. A tensão nominal é 138 quilovolts. Em cada local observamos as fases A, B e C. Cada execução fornece nove máximos, com seus instantes e identificadores. Esses registros compartilham a mesma execução; não são nove simulações independentes. São resumos dos picos, não formas de onda completas. A base reúne dois cenários desse sistema.
\cue{Explicar o ponto central e fazer a transição ao próximo slide.}


\slide{6}{Cenários de simulação e composição da base}{3:40--4:25}
CRPI significa com resistor de pré-inserção, usado durante a conexão para atenuar o transitório. SRPI significa sem esse componente. Cada cenário possui arquivos com 50, 100, 200, mil e dez mil simulações. São dez arquivos e 204.300 máximos, organizados em 90 grupos no total: dois cenários, cinco tamanhos, três locais e três fases. Os cenários também diferem no chaveamento; a comparação não isola o efeito causal do resistor. Para analisar as amplitudes, usamos uma referência comum.
\cue{Explicar o ponto central e fazer a transição ao próximo slide.}


\slide{7}{Normalização e interpretação dos valores de tensão}{4:25--5:15}
Por unidade, ou P.U., é o valor dividido por sua referência. Um resultado de 1,0 coincide com ela; 1,5 fica 50 por cento acima. A base de pico fase-terra é 112.677 volts. Um pico ilustrativo de 225.354 volts corresponde a 2,0 P.U. Os 138 quilovolts nominais são uma grandeza diferente: tensão eficaz entre fases. Mantemos também sinal original, instante, local, fase e execução. Essa transformação integra o fluxo computacional mostrado a seguir.
\cue{Explicar o ponto central e fazer a transição ao próximo slide.}


\slide{8}{Tecnologias e fluxo de processamento}{5:15--6:10}
A implementação usa Python para integrar o processamento. Polars organiza os registros em tabelas; NumPy e SciPy apoiam cálculos numéricos e diagnósticos estatísticos. scikit-learn fornece K-Means, DBSCAN e índices dos agrupamentos. A seleção hierárquica da cauda superior é uma regra própria em Python, baseada nas lacunas entre valores ordenados; não usa o agrupamento de Ward do scikit-learn. Matplotlib produz gráficos, e Pytest verifica regras do software.

O fluxo parte dos arquivos LIS, extrai e valida registros, converte amplitudes para P.U. e produz estatísticas, agrupamentos, tabelas e gráficos. A identidade de origem acompanha cada etapa. As bibliotecas dão suporte; a contribuição está na integração e nas verificações específicas dos dados do ATP. Com essa base pronta, separamos as duas perguntas analíticas.
\cue{Percorrer o fluxo; associar cada tecnologia a uma função, sem ler uma lista de ferramentas.}


\slide{9}{Excedência e discrepância: objetivos distintos}{6:10--7:00}
O fluxo responde a duas perguntas. Excedência significa ultrapassar o limiar fixo de 2,3 P.U.; contamos eventos por grupo. Discrepância significa destoar dos demais valores no mesmo contexto. No exemplo ilustrativo, 1,8 P.U. pode destoar de um grupo próximo de 1,2 sem superar 2,3. Candidato não significa erro confirmado, e o limiar não recomenda isolamento. Começaremos pela distribuição dos máximos e pela estimativa de excedência; depois veremos agrupamentos e candidatos.
\cue{Explicar o ponto central e fazer a transição ao próximo slide.}


\slide{10}{Distribuição dos máximos e adequação do modelo normal}{7:00--8:25}
Na análise probabilística, vamos olhar primeiro o gráfico da esquerda. O eixo horizontal mostra os
máximos em P.U.; o eixo vertical mostra densidade, uma frequência
normalizada. A área das barras representa a proporção em cada faixa.
Não são contagens brutas. O gráfico mostra onde os máximos se concentram.

A curva sobreposta tem a conhecida forma de sino da distribuição normal.
Ela foi ajustada para comparação, mas neste grupo não acompanha bem
os dados. A referência de Hevia fundamenta a interpretação das saídas
do ATP; não garante que os máximos sigam uma distribuição normal. O painel da direita é outra comparação com a normal:
se o ajuste fosse adequado, os pontos tenderiam a acompanhar a reta.

O diagnóstico estatístico rejeitou a hipótese de normalidade em 88 dos
90 grupos, ao nível de cinco por cento. As duas não rejeições, com
amostras pequenas, não provam normalidade. A conclusão também não é
que todos os dados sejam anômalos.

Média e desvio padrão continuam úteis para descrever os valores. Mas,
para estimar a chance de ultrapassar o limiar, o resultado principal usa
contagens observadas e um intervalo de incerteza. Não impomos a curva
normal como se ela explicasse os dados nem a usamos para escolher
candidatos.
\cue{Apontar eixos e barras do histograma. Só depois explicar brevemente o painel à direita.}


\slide{11}{Probabilidade de excedência e incerteza amostral}{8:25--9:25}
Retomando a primeira pergunta: algum máximo passou de 2,3 P.U.?
Não houve nenhuma ocorrência nos 90 grupos. A frequência observada foi
zero, mas isso não demonstra que o evento seja impossível.

É como executar testes de software sem encontrar uma falha: isso informa
sobre os testes realizados, mas não prova que nenhuma falha possa existir.
A analogia serve para lembrar o limite da observação.

Para expressar a incerteza, calculamos um intervalo bilateral de Wilson de
95 por cento em cada grupo. Com zero eventos em 50 simulações, o limite
superior fica em aproximadamente 7,13 por cento. Com zero em dez mil,
cai para aproximadamente 0,04 por cento.

Esses números não são probabilidades observadas; são limites do intervalo
estimado pelo procedimento. O resultado central é que uma amostra maior
estreita esse intervalo, sob as condições do modelo estatístico.
Não reunimos as 204.300 observações em um único denominador: cada grupo
mantém seu número de simulações e seu intervalo. Respondida a pergunta de excedência, passamos à discrepância e aos métodos de agrupamento.
\cue{Comparar as duas linhas: mesmo zero observado, incertezas diferentes.}


\slide{12}{Métodos de clusterização avaliados}{9:25--11:00}
Passamos agora à comparação dos métodos. K-Means divide os valores em dois agrupamentos. A candidatura usa outra regra: distância ao centroide estritamente acima do percentil 99 das distâncias de calibração. O grupo de maior tensão não é automaticamente anômalo, e a regra não estima a proporção de erros reais.

DBSCAN agrupa por densidade e sinaliza ruído quando o ajuste é aplicável. O pacote usa raio de 0,5 no espaço padronizado e mínimo de cinco vizinhos. A regra hierárquica publicada na monografia separa até dez por cento da cauda superior, conservando três pontos de referência. A lacuna deve superar 0,05 P.U. e cinco vezes a mediana das lacunas.

Os três métodos têm finalidades diferentes. Vamos observar suas partições sobre os mesmos dados antes de comparar os índices.
\cue{Distinguir rótulo de agrupamento e regra de candidatura. Parâmetros detalhados ficam disponíveis para perguntas.}


\slide{13}{Agrupamentos: K-Means}{11:00--12:10}
Este gráfico usa CRPI, fase A e 200 simulações. O eixo horizontal representa a posição relativa na linha: zero, meio e terminal oposto; o vertical mostra máximos em P.U. Cada terminal é ajustado separadamente. A posição serve à visualização, não é atributo do agrupamento. O K-Means produz dois grupos por terminal. O grupo de maior tensão não é automaticamente anômalo: candidatos seguem a distância ao centroide acima do percentil 99. No próximo gráfico, observamos os mesmos dados por densidade.
\cue{Apontar os grupos no terminal oposto; manter a comparação com o gráfico anterior.}

\slide{14}{Agrupamentos: DBSCAN}{12:10--13:20}
Mantemos cenário, fase, tamanho e escalas. O DBSCAN reúne a região densa em um agrupamento e marca pontos de ruído, visíveis no terminal oposto. Ruído não significa erro confirmado. Mesmo com candidatos, a partição pode não permitir Silhouette e Davies--Bouldin: após retirar o ruído, resta um único agrupamento. Essa distinção explica por que ausência de índice não equivale a falha do detector. Agora veremos a regra hierárquica sobre os mesmos dados.
\cue{Apontar os grupos no terminal oposto; manter a comparação com o gráfico anterior.}

\slide{15}{Agrupamentos: Hierárquico}{13:20--14:25}
A seleção hierárquica procura uma pequena cauda superior separada por uma lacuna suficiente. Neste recorte, ela aparece no terminal oposto, enquanto os demais valores permanecem no grupo de referência. A implementação usa ligação simples em uma dimensão, pela análise das lacunas entre valores ordenados. As cores são específicas deste método; não correspondem às mesmas classes do K-Means. Os três gráficos mostram finalidades diferentes sobre dados iguais. Para quantificar essas diferenças, passamos aos índices geométricos.
\cue{Apontar os grupos no terminal oposto; manter a comparação com o gráfico anterior.}

\slide{16}{Métricas principais: Silhouette e Davies--Bouldin}{14:25--15:45}
As métricas principais da comparação são Silhouette e Davies--Bouldin.
Silhouette compara a proximidade de cada ponto ao próprio grupo com a
proximidade ao grupo vizinho. Valores próximos de um indicam melhor
coesão e separação; valores próximos de zero sugerem fronteiras pouco
nítidas, e negativos indicam que outro grupo pode ser mais próximo.

Davies--Bouldin relaciona a dispersão interna à separação entre centros.
Nesse índice, valores menores são preferíveis. As métricas observam
aspectos diferentes, portanto podem produzir ordenações distintas.

Antes de comparar os números, conferimos se a partição permite o cálculo.
Na política adotada, um único grupo ou a presença de grupo unitário
impede a avaliação. No DBSCAN, os pontos de ruído são retirados desse
cálculo. Um índice não aplicável não recebe zero e não significa falha
de execução. Nenhuma das métricas mede probabilidade de erro elétrico.
\cue{Fixar direções: Silhouette maior; Davies--Bouldin menor.}


\slide{17}{Comparação dos índices em recortes específicos}{15:45--17:05}
A primeira tabela compara K-Means e hierárquico no cenário com resistor,
terminal oposto, fase A, com 200 observações. Ambos
avaliam exatamente os mesmos valores.
O K-Means tem Silhouette de 0,7127, acima de 0,5444 do hierárquico.
Já Davies--Bouldin favorece o hierárquico: 0,3182 contra 0,4065.
Portanto, os índices não apontam necessariamente o mesmo vencedor.

A segunda tabela usa outro recorte: com resistor, meio da linha,
fase A, com 50 simulações. K-Means e DBSCAN avaliam os mesmos 50 valores,
sem exclusão de ruído pelo DBSCAN. Silhouette favorece o K-Means:
0,7736 contra 0,7580. Davies--Bouldin favorece o DBSCAN:
0,2381 contra 0,3928.

As tabelas mostram somente pares com os dois índices aplicáveis.
Isso facilita a leitura; a abrangência da avaliação dos três métodos
será apresentada no próximo slide.

Não devemos comparar diretamente as duas tabelas como se representassem
o mesmo experimento. Em cada tabela, a população é a mesma entre os
métodos comparados, mas os índices favorecem métodos diferentes.
Os resultados não sustentam um vencedor universal.
\cue{Setas indicam direção preferível; negrito destaca melhor índice. Comparar dentro de cada tabela e apontar denominadores.}


\slide{18}{Tamanho amostral e abrangência da comparação}{17:05--18:10}
Esta tabela examina o tamanho amostral no cenário sem resistor para o
K-Means. Cada linha mostra as medianas dos nove grupos terminal--fase,
não o resultado de um único agrupamento de todas as observações.
Estão presentes os cinco tamanhos: 50, 100, 200, mil e dez mil.

Silhouette passa de 0,6491 com 50 simulações para 0,6770 com cem e
0,6258 com dez mil. Davies--Bouldin passa de 0,4499 para 0,4663 e
0,5032. Portanto, aumentar a amostra não melhora automaticamente a
separação dos grupos. No K-Means, o número de agrupamentos continua fixo
em dois.

O painel da direita mostra outra dimensão: a abrangência da avaliação. Os índices foram
aplicáveis aos 90 grupos do K-Means, a quatro do DBSCAN e a seis do
hierárquico. Isso limita comparações globais: uma mediana de poucos
grupos não representa toda a base com a mesma cobertura.

Ausência de índice não significa que o método falhou. Pode significar
que produziu uma partição sem estrutura elegível para essas métricas.
A comparação entre tamanhos é descritiva; não demonstra efeito causal
isolado da quantidade de simulações.
\cue{Distinguir qualidade da partição, abrangência dos índices e tamanho amostral.}

Essa comparação descreve geometria e sua abrangência. A seguir, reunimos as contribuições computacionais e um caso rastreável.


\slide{19}{Contribuições computacionais e delimitação do escopo}{18:10--19:15}
A rotina integrou extração, validação, normalização e análise nos dez arquivos, preservando 204.300 observações. Sessenta valores individuais da fase A e noventa médias concordaram com ATP; isso não representa auditoria individual de toda a base.

A união dos três critérios publicados no pacote v10 reuniu 2.010 candidatos exploratórios, sem causas confirmadas. Um exemplo é a execução 43 da fonte SRPI com 50 simulações, terminal oposto, fase C: 1,2467 P.U., sinalizado pelo DBSCAN e abaixo de 2,3 P.U. Recuperamos sua origem e o motivo da sinalização.

A entrega é uma base verificável para investigação. Não houve exclusão automática por candidatura nem recomendação de isolamento. Esses resultados permitem concluir sobre o alcance da comparação.
\cue{Apresentar o total como resultado do pacote publicado; não como quantidade de erros.}


\slide{20}{Conclusões e continuidade}{19:15--20:00}
O K-Means apresentou maior cobertura dos índices, mas isso não demonstra superioridade na detecção. Silhouette e Davies--Bouldin podem favorecer métodos diferentes. As finalidades e populações avaliadas impedem declarar um vencedor universal.

A contribuição é uma rotina que seleciona casos e preserva evidências. A continuidade exige formas de onda, repetições com menor passo de integração e revisão elétrica independente. Os casos permanecem sem conclusão causal. Obrigado; fico à disposição para perguntas.
\cue{Encerrar com resultado comparativo, contribuição e limite.}

\section*{Apoio para perguntas da banca}

As respostas abaixo são apoio ao estudo, fora do tempo de exposição.

\subsection*{O trabalho detectou erros numéricos reais do ATP?}
Não confirmou causas. Produziu candidatos analíticos rastreáveis.
A confirmação exige formas de onda, verificação das condições de
chaveamento, repetições comparáveis e revisão elétrica qualificada.
Os casos de revisão permanecem com conclusão causal não resolvida.

\subsection*{Qual é a contribuição acadêmica se a causa não foi confirmada?}
A contribuição está na rotina integrada de extração, validação e análise,
na preservação da identidade das observações, na estimativa empírica
com incerteza e na comparação explícita dos critérios. Ela prepara uma
base auditável para investigação elétrica posterior, dentro do escopo
computacional do TCC.

\subsection*{Por que a base é 112.677 V e não 138.000 V?}
138 kV é a tensão nominal eficaz entre fases. Para máximos fase-terra,
a referência de pico é $138.000\sqrt{2/3}$ volts, aproximadamente
112.676,528 V. O ATP imprime a base arredondada de 112.677 V, utilizada
para reproduzir seus resumos. A comparação de 60 pontos e 90 médias
sustenta a conversão no escopo auditado.

\subsection*{Normalidade rejeitada invalida média e desvio padrão?}
Não. Essas medidas permanecem descritivas. A rejeição limita o uso da
normal ajustada para inferência probabilística. Não demonstra ausência
de qualquer distribuição, nem confirma anomalia ou erro em uma observação.

\subsection*{Se não houve excedência de 2,3 P.U., por que existem candidatos?}
Excedência e discrepância são perguntas diferentes. A primeira compara
o máximo com um limiar fixo. Os critérios de seleção avaliam distância
ao centroide, ruído por densidade ou pertença à cauda superior hierárquica.
Um ponto pode
ser discrepante dentro de seu grupo e permanecer abaixo de 2,3 P.U.

\subsection*{Por que não usar três desvios padrão para detectar tensões?}
No modelo examinado, o parâmetro de três desvios padrão pertence ao
chaveamento estatístico de entrada. Ele não define um critério das tensões
de saída. No pacote apresentado, a seleção combina distância do
K-Means, ruído do DBSCAN e pertença à cauda superior hierárquica,
sob suas regras de aplicabilidade. Nenhum desses critérios confirma
isoladamente uma causa elétrica ou numérica.

\subsection*{O K-Means é melhor por ter índices em todos os grupos?}
Não é possível concluir isso. Ele fixa duas partições na comparação,
enquanto DBSCAN pode produzir um único agrupamento e ruído. A
aplicabilidade dos índices mede se a partição permite o cálculo,
sem avaliar diretamente a acurácia da seleção. O desempenho sintético
também se restringe às condições controladas testadas.

\subsection*{O percentil 99 do K-Means implica cerca de 1\% de erros?}
Não. O limiar estabelece uma regra de seleção relativa às distâncias
de calibração. A fração selecionada depende da comparação estrita,
de empates e da distribuição das distâncias. Ela não estima a
prevalência de erros numéricos. A seleção exploratória usa os mesmos
dados para calibração e ajuste; não constitui teste independente.

\subsection*{Jaccard igual a um comprova que o método está correto?}
Comprova apenas que os conjuntos selecionados coincidiram na grade
testada e no recorte avaliado. Um conjunto pode ser estável e ainda
não corresponder a erros reais. É evidência de estabilidade da seleção,
sem validação causal.

\subsection*{O resistor reduziu as tensões?}
No recorte apresentado, a média do CRPI foi menor que a do SRPI.
Entretanto, os cenários preservam definições distintas de chaveamento.
O resultado descreve a comparação observada e não isola o efeito causal
do resistor por um experimento pareado que controle todos os demais fatores.

\subsection*{Como reproduzir e quais limites permanecem?}
O pacote de evidências registra fontes, parâmetros e artefatos em um manifesto de
reprodutibilidade. A aceitação vale para os dez arquivos, seus leiautes
e pontos de verificação declarados. Novos formatos e condições exigem
nova validação. O estado do pacote continua incompleto quanto à
validação elétrica dos candidatos, conforme a monografia.

\section*{Orientação para o ensaio}
Cronometre com os slides abertos. Marcos: 6:10 ao terminar tecnologias (slide 8), 9:25 ao terminar probabilidade (slide 11), 14:25 ao terminar os três gráficos (slide 15) e 18:10 ao terminar abrangência (slide 18). Encerramento começa em 19:15 (slide 20). A tabela inicial contém todos os tempos.

Se ficar longa, reduza analogias iniciais e leitura de números. Preserve os três gráficos, a explicação dos critérios, a comparação dos índices e a conclusão. Se ficar curta, explore as diferenças visíveis entre as partições e suas implicações para a aplicabilidade dos índices.

\section*{Detalhes técnicos reservados para perguntas}
A rotina usa Python, Polars, SciPy e scikit-learn. A validação verifica
ausências, duplicidades e valores inválidos antes dos modelos.
As 204.300 observações foram preservadas, sem exclusão automática.
Na conferência numérica, 60 valores individuais da fase A e 90 médias
concordaram com o ATP; isso não equivale à auditoria individual de toda
a base.

O K-Means usa dois
agrupamentos e distância ao centroide acima do percentil 99 da calibração.
O DBSCAN usa $\varepsilon=0{,}5$ no espaço padronizado e mínimo de cinco
vizinhos, sob regras de aplicabilidade. O hierárquico seleciona os membros
de uma pequena cauda superior separada por uma lacuna suficiente.
Esse critério participa da união de candidatos; os índices geométricos
avaliam a partição em uma análise distinta.

Os índices Silhouette e Davies--Bouldin avaliam separação geométrica,
não causas elétricas. Foram aplicáveis a 90 grupos do K-Means, quatro do
DBSCAN e seis da seleção hierárquica de cauda superior. Ausência de índice
não equivale a detector não aplicável. Vinte réplicas na grade de
sensibilidade tiveram Jaccard igual a um no recorte testado, evidência
de estabilidade que não comprova causa.

A taxa de candidatos no CRPI, outra ponta, fase A, foi de 0\% em
50 simulações e 0,99\% em 10.000, sem evolução monotônica entre tamanhos. As fontes não são presumidas independentes
ou aninhadas. A amostra maior é referência empírica, não verdade fundamental.
O modelo elétrico é K.C. Lee não transposto, com passo de integração de
$10\,\mu$s e janela simulada de $0{,}3$ s.

\section*{Distinção entre avaliação e seleção}
Silhouette e Davies--Bouldin são as métricas de avaliação dos agrupamentos
utilizadas na apresentação. Medem coesão e separação das partições; não
selecionam individualmente candidatos nem confirmam suas causas.
A comparação geométrica é distinta da comparação entre as sinalizações
dos métodos.

O hierárquico participa da seleção de candidatos pela cauda superior,
conforme a metodologia da monografia e o pacote v10. Essa participação
não equivale à confirmação de erro numérico. O exemplo individual do
slide 19 utiliza a sinalização explícita do DBSCAN.

\section*{Fontes para preparação}
A apresentação se baseia na monografia \texttt{doc/main.tex} e no pacote
de evidências, especialmente as tabelas de estatística, probabilidade,
comparação de métodos e ensaios controlados. Consultar as referências
bibliográficas da monografia para os fundamentos dos métodos.

\section*{Compilação}
A partir do diretório \texttt{doc/}, execute:
\begin{verbatim}
latexmk -pdf -outdir=build/apresentacao apresentacao.tex
latexmk -pdf -outdir=build/apresentacao roteiro_apresentacao.tex
\end{verbatim}
Os PDFs serão gerados em \texttt{doc/build/apresentacao/}.

\clearpage
\section*{Alternativa: tópicos para apresentar com suas próprias palavras}
Use os tópicos na mesma ordem e nos mesmos tempos do roteiro.

\slide{1}{Abertura}{0:00--0:30}
\begin{itemize}
\item Apresentação pessoal, orientação e coorientação.
\item Tema: processamento de simulações elétricas e análise de anomalias.
\item Percurso da apresentação: contexto, solução computacional e resultados.
\end{itemize}

\slide{2}{Contextualização: energização de linhas de transmissão}{0:30--1:15}
\begin{itemize}
\item Linha de transmissão: transporte de energia entre subestações.
\item Energização: conexão da linha à rede pelo fechamento de uma chave.
\item Transitório: oscilações logo após a conexão; definir sobretensão.
\item Relevância: isolamento precisa suportar esses esforços elétricos.
\end{itemize}

\slide{3}{Simulação elétrica e problema computacional}{1:15--2:05}
\begin{itemize}
\item Ensaios reais: custos e restrições operacionais.
\item ATP simula transitórios; ATPDraw representa o circuito graficamente.
\item Repetições com diferentes instantes de chaveamento geram arquivos LIS.
\item Desafio: extrair, validar e comparar dados, preservando a origem.
\item Caso discrepante exige investigação; o número isolado não revela a causa.
\end{itemize}

\slide{4}{Objetivos específicos}{2:05--2:50}
\begin{itemize}
\item Extrair e organizar máximos, preservando execução, local e fase.
\item Pré-processar; média, desvio padrão e dispersão dos máximos.
\item Estimar frequências de excedência e intervalos de confiança.
\item Aplicar K-Means, DBSCAN e hierárquico: grupos e candidatos atípicos.
\item Avaliar agrupamentos com Silhouette e Davies--Bouldin;
comparar resultados estatísticos e sinalizações dos métodos.
\item Integrar análise e apresentação dos resultados em rotina reprodutível.
\item Transição: distinguir excedência de discrepância no próprio grupo.
\end{itemize}

\slide{5}{Estudo de caso e pontos de observação}{2:50--3:40}
\begin{itemize}
\item Apontar no esquema: terminal de manobra, meio e terminal oposto.
\item Linha de 80 km, tensão nominal de 138 kV e fases A, B e C.
\item Nove máximos por execução: três locais multiplicados por três fases.
\item Máximo e instante resumem o pico; não descrevem a forma de onda completa.
\item Nove máximos da mesma execução não são nove execuções independentes.
\end{itemize}

\slide{6}{Cenários de simulação e composição da base}{3:40--4:25}
\begin{itemize}
\item Definir CRPI e SRPI: com e sem resistor de pré-inserção.
\item Por cenário: 50, 100, 200, 1.000 e 10.000 simulações.
\item Dez arquivos, 204.300 máximos e 90 grupos de análise ao todo.
\item Separação por cenário, tamanho amostral, local e fase.
\item Chaveamento também difere: comparação não isola o efeito do resistor.
\end{itemize}

\slide{7}{Normalização e interpretação dos valores de tensão}{4:25--5:15}
\begin{itemize}
\item P.U.: valor dividido pela referência; facilita comparar amplitudes.
\item Exemplos: 1,0 é a referência; 1,5 é 50\% acima; 2,0 é o dobro.
\item Base de pico fase-terra: 112.677 V; distinguir dos 138 kV entre fases.
\item Exemplo ilustrativo: 225.354 V correspondem a 2,0 P.U.
\item Preservar instante, local, fase e execução junto ao máximo.
\end{itemize}

\slide{8}{Tecnologias e fluxo de processamento}{5:15--6:10}
\begin{itemize}
\item Fluxo: LIS, extração, validação, P.U., análise e artefatos.
\item Tecnologias associadas a funções; identidade preservada.
\item Hierárquico: regra própria em Python; K-Means e DBSCAN: scikit-learn.
\end{itemize}

\slide{9}{Excedência e discrepância: objetivos distintos}{6:10--7:00}
\begin{itemize}
\item Excedência: ultrapassar o limiar fixo de 2,3 P.U.
\item Discrepância: destoar dos demais registros do próprio grupo.
\item Exemplo conceitual: 1,8 P.U. pode destoar de valores próximos de
1,2 P.U. sem ultrapassar 2,3 P.U.
\item Candidato não confirma erro; limiar não recomenda nível de isolamento.
\end{itemize}

\slide{10}{Distribuição dos máximos e adequação do modelo normal}{7:00--8:25}
\begin{itemize}
\item Ler o histograma: máximos no eixo horizontal, densidade nas barras;
área representa proporção, não contagem bruta.
\item Curva normal e painel à direita: comparações de adequação do modelo.
\item Referência de Hevia orienta a interpretação das saídas do ATP;
não garante normalidade.
\item Normalidade rejeitada em 88 de 90 grupos a 5\%; duas não rejeições
não provam normalidade.
\item Média e desvio padrão permanecem descritivos; excedência principal
estimada por frequência empírica, sem detector baseado na normal.
\end{itemize}

\slide{11}{Probabilidade de excedência e incerteza amostral}{8:25--9:25}
\begin{itemize}
\item Zero máximos acima de 2,3 P.U. nos 90 grupos.
\item Zero observado não demonstra impossibilidade do evento.
\item Wilson de 95\%: limite superior de 7,13\% com 50 simulações e
aproximadamente 0,04\% com 10.000, ambos com zero eventos.
\item Mesma frequência observada, menor incerteza com amostra maior.
\item Calcular por grupo; não usar 204.300 como denominador único.
\end{itemize}

\slide{12}{Métodos de clusterização avaliados}{9:25--11:00}
\begin{itemize}
\item K-Means: dois agrupamentos por proximidade aos centros.
\item DBSCAN: agrupamentos por densidade e identificação de ruído.
\item Hierárquico: ligação simples em uma dimensão; membros da cauda
superior separada são candidatos, sob regras de aplicabilidade.
\item Comparar partições por Silhouette e Davies--Bouldin.
\item Boa separação geométrica não comprova erro numérico ou causa elétrica.
\item K-Means: candidato pela distância acima do percentil 99, não pelo rótulo.
\end{itemize}

\slide{13}{Agrupamentos: K-Means}{11:00--12:10}
\begin{itemize}
\item Este gráfico usa CRPI, fase A e 200 simulações. O eixo horizontal representa a posição relativa na linha: zero, meio e terminal oposto; o vertical mostra máximos em P.U. Cada terminal é ajustado separadamente. A posição serve à visualização, não é atributo do agrupamento. O K-Means produz dois grupos por terminal. O grupo de maior tensão não é automaticamente anômalo: candidatos seguem a distância ao centroide acima do percentil 99. No próximo gráfico, observamos os mesmos dados por densidade.
\end{itemize}

\slide{14}{Agrupamentos: DBSCAN}{12:10--13:20}
\begin{itemize}
\item Mantemos cenário, fase, tamanho e escalas. O DBSCAN reúne a região densa em um agrupamento e marca pontos de ruído, visíveis no terminal oposto. Ruído não significa erro confirmado. Mesmo com candidatos, a partição pode não permitir Silhouette e Davies--Bouldin: após retirar o ruído, resta um único agrupamento. Essa distinção explica por que ausência de índice não equivale a falha do detector. Agora veremos a regra hierárquica sobre os mesmos dados.
\end{itemize}

\slide{15}{Agrupamentos: Hierárquico}{13:20--14:25}
\begin{itemize}
\item A seleção hierárquica procura uma pequena cauda superior separada por uma lacuna suficiente. Neste recorte, ela aparece no terminal oposto, enquanto os demais valores permanecem no grupo de referência. A implementação usa ligação simples em uma dimensão, pela análise das lacunas entre valores ordenados. As cores são específicas deste método; não correspondem às mesmas classes do K-Means. Os três gráficos mostram finalidades diferentes sobre dados iguais. Para quantificar essas diferenças, passamos aos índices geométricos.
\end{itemize}

\slide{16}{Métricas principais: Silhouette e Davies--Bouldin}{14:25--15:45}
\begin{itemize}
\item Silhouette: coesão e separação; quanto maior, melhor.
\item Davies--Bouldin: dispersão relativa à separação; quanto menor, melhor.
\item Métricas podem favorecer partições diferentes.
\item Política adotada: um grupo ou grupo unitário impede o cálculo;
ruído do DBSCAN fica fora da população avaliada.
\item Não aplicável não significa zero, falha de execução ou erro elétrico.
\end{itemize}

\slide{17}{Comparação dos índices em recortes específicos}{15:45--17:05}
\begin{itemize}
\item Primeira tabela: mesmo recorte, índices favorecem métodos diferentes
entre K-Means e hierárquico.
\item Segunda tabela: com resistor, meio da linha, fase A, 50 simulações;
K-Means e DBSCAN avaliam os mesmos 50 valores, sem exclusão de ruído.
\item Silhouette favorece K-Means; Davies--Bouldin favorece DBSCAN.
\item Mostrar somente pares com índices aplicáveis; cobertura dos três
métodos será discutida no próximo slide.
\item Comparar dentro de cada tabela, conferir denominadores;
não há vencedor universal.
\end{itemize}

\slide{18}{Tamanho amostral e abrangência da comparação}{17:05--18:10}
\begin{itemize}
\item Esquerda: medianas dos nove grupos terminal--fase, K-Means sem resistor;
cinco tamanhos amostrais, sem omissões.
\item Mais simulações não melhoram automaticamente a separação geométrica.
\item Cobertura dos índices: K-Means 90/90, DBSCAN 4/90, hierárquico 6/90.
\item Medianas de poucos grupos não representam a base inteira.
\item Comparação descritiva entre fontes; não isola efeito causal do tamanho
amostral nem demonstra convergência por si só.
\end{itemize}

\slide{19}{Contribuições computacionais e delimitação do escopo}{18:10--19:15}
\begin{itemize}
\item Dez arquivos, 204.300 observações preservadas.
\item Concordância com ATP: 60 valores da fase A e 90 médias; cobertura parcial.
\item Exemplo: SRPI, 50 simulações, execução 43, terminal oposto, fase C;
1,2467 P.U., sinalização DBSCAN, sem exceder 2,3 P.U.
\item Intervalos por grupo e comparação geométrica com cobertura explícita.
\item Contribuição: apoiar investigação dos resultados de simulação.
\item Escopo exclui recomendação de isolamento e exclusão automática de dados.
\item União publicada no pacote v10: 2.010 candidatos; não são erros confirmados.
\end{itemize}

\slide{20}{Conclusões e continuidade}{19:15--20:00}
\begin{itemize}
\item Maior cobertura do K-Means não prova melhor detecção.
\item Índices discordam; não há vencedor universal.
\item Evidências rastreáveis; revisão causal continua pendente.
\end{itemize}

\end{document}
